\documentclass{beamer}
\usepackage{graphicx}
\usepackage{xcolor}
\usepackage{tikz}

% 自定义颜色
\definecolor{purpleblue}{RGB}{80, 80, 180}
\definecolor{lightgray}{RGB}{230, 230, 230}
\definecolor{novelPurple}{RGB}{98, 54, 255} 
\definecolor{darkgray}{RGB}{50, 50, 50}

% 主题设置
\usetheme{Madrid}  
\usecolortheme{dolphin} 

% 设置字体（移除 fontspec，确保 pdfLaTeX 兼容）
\setbeamerfont{title}{size=\Large,series=\bfseries}
\setbeamerfont{frametitle}{size=\LARGE,series=\bfseries}
\setbeamercolor{frametitle}{fg=white, bg=novelPurple}
\setbeamercolor{title}{fg=white, bg=purpleblue}

% 封面信息
\title[SAT Hard Question]{\textbf{SAT Hard Question 13}}
\author{\textbf{Jack Zeng}}
\institute{\textcolor{white}{\textbf{Novel Prep}}}
\date{\today} 

\begin{document}

% --- Title Slide ---
\begin{frame}
    \titlepage
    \vfill
    \begin{flushright}
        \textcolor{purpleblue}{\Huge \textbf{Novel Prep}}
    \end{flushright}
\end{frame}


\begin{frame}
    \begin{tikzpicture}[remember picture, overlay]
        % 设置浅色渐变背景
        \shade[top color=softPurple, bottom color=softGray] 
            (current page.north west) rectangle (current page.south east);
        
        % 添加高斯模糊光晕
        \fill[white, opacity=0.3] (current page.center) circle (4cm);
        
        % 主标题
        \node[align=center] at (current page.center) {
            {\Huge\bfseries\textcolor{purpleblue}{Novel Prep}} \\[0.5cm]
            {\Large\itshape\textcolor{lightText}{Excellence in SAT Prep}}
        };

        % 添加柔和光点（点缀）
        \foreach \x in {1,2,3,4,5} {
            \fill[purpleblue, opacity=0.2] (rand*10-5, rand*7-3.5) circle (0.1+\x*0.05);
        }

        ;
    \end{tikzpicture}
\end{frame}


\begin{frame}{\textbf{Question: Finding Circumference from Arc Length}}
A circle has center \( O \), and points \( A \) and \( B \) lie on the circle. The measure of arc \( AB \) is \( 60^\circ \), and the length of this arc is 4 inches.  

What is the circumference, in inches, of the circle?

\begin{enumerate}[A.]
    \item 4
    \item 8
    \item 16
    \item 24
\end{enumerate}
\end{frame}


\begin{frame}{\textbf{Answer Explanation: Finding Circumference from Arc Length (Part 1)}}
\textbf{Step 1:}  
The formula for arc length is:  
\[
\text{Arc Length} = \frac{\theta}{360^\circ} \cdot \text{Circumference}
\]  
where \( \theta = 60^\circ \) and arc length is 4 inches.  
Let \( C \) be the circumference.

\pause

\textbf{Step 2:} Plug into the formula:  
\[
4 = \frac{60}{360} \cdot C
\Rightarrow 4 = \frac{1}{6}C
\Rightarrow C = 4 \cdot 6 = \boxed{24}
\]

\textbf{Correct Answer:} \textbf{(D) 24}
\end{frame}



\begin{frame}{\textbf{Question: X-Intercepts of a Rational Function}}
In the xy-plane, the graph of function \( f \), where \( y = f(x) \), has exactly 8 x-intercepts. One of these x-intercepts is \( (17, 0) \).

The rational function \( g \) is defined by:  
\[
g(x) = \frac{f(x)}{x - 17}
\]

In the xy-plane, how many x-intercepts does the graph of \( y = g(x) \) have?

\begin{enumerate}[A.]
    \item 5
    \item 6
    \item 7
    \item 8
\end{enumerate}
\end{frame}


\begin{frame}{\textbf{Answer Explanation: X-Intercepts of Rational Functions (Part 1)}}
\textbf{Step 1:}  
The original function \( f(x) \) has 8 x-intercepts, including \( x = 17 \).  
This means \( f(17) = 0 \), so 17 is a zero of \( f(x) \).  

\pause

\textbf{Step 2:}  
In the function \( g(x) = \frac{f(x)}{x - 17} \), the zero at \( x = 17 \) is now in both the numerator and the denominator.  
This creates a removable discontinuity (a hole), not an x-intercept.

\textbf{Step 3:}  
The point \( (17, 0) \) is no longer an x-intercept of \( g(x) \).  
So, \( g(x) \) has \( 8 - 1 = \boxed{7} \) x-intercepts.

\textbf{Correct Answer:} \textbf{(C) 7}
\end{frame}


\begin{frame}{\textbf{Question: Predator and Prey Growth Percent}}
A researcher investigated two species of mites: a predator and its prey. At the start of a week, there was an equal number of the two species.

At the end of the week:
\begin{itemize}
    \item The number of prey had increased by \( 2700\% \) of the number of prey at the start of the week.
    \item The number of predators had increased by \( 180\% \) of the number of predators at the start of the week.
\end{itemize}

The number of prey at the end of the week was \( p\% \) greater than the number of predators at the end of the week.

What is the value of \( p \)?

\begin{enumerate}[A.]
    \item 9
    \item 90
    \item 900
    \item 9000
\end{enumerate}
\end{frame}

\begin{frame}{\textbf{Answer Explanation: Predator and Prey Growth Percent (Part 1)}}
\textbf{Step 1: Assume initial populations.}

Let the initial number of prey and predators be \( x \).  
\pause

\textbf{Step 2: Compute population at the end of the week.}
\begin{itemize}
    \item Prey: \( x + 2700\% \cdot x = x + 27x = 28x \)
    \item Predators: \( x + 180\% \cdot x = x + 1.8x = 2.8x \)
\end{itemize}
\pause

\textbf{Step 3: Use percentage increase formula.}

The prey population is \( p\% \) greater than the predator population:
\[
\frac{28x - 2.8x}{2.8x} = \frac{25.2x}{2.8x} = 9
\Rightarrow p = 900
\]

\textbf{Correct Answer:} \textbf{(C) 900}
\end{frame}




\begin{frame}{\textbf{Question: Solving for Base of an Exponential Function}}
The exponential function \( f \) is defined by
\[
f(x) = ab^x,
\]
where \( a \) and \( b \) are positive constants.

If \( f(s) = t \) and \( f(s+1) = t - 0.87t \), where \( s \) and \( t \) are constants, what is the value of \( b \)?

\begin{enumerate}[A.]
    \item 0.13
    \item 0.87
    \item 1.13
    \item 1.87
\end{enumerate}
\end{frame}



\begin{frame}{\textbf{Answer Explanation: Exponential Function (Part 1)}}
We are given:
\[
f(s) = ab^s = t, \quad f(s+1) = ab^{s+1} = t - 0.87t = 0.13t
\]
\pause

Now divide the two expressions:
\[
\frac{f(s+1)}{f(s)} = \frac{ab^{s+1}}{ab^s} = \frac{0.13t}{t}
\Rightarrow \frac{b^{s+1}}{b^s} = b = 0.13
\]

\textbf{Correct answer: (A) 0.13}
\end{frame}


\begin{frame}{\textbf{Question: Bee Population Prediction Model}}
A beekeeper’s initial observation of the population of a certain bee colony was 1,100 bees. The beekeeper set a goal to increase the population to 2,600 bees.

The beekeeper uses a model that predicts the population of this bee colony:
\begin{itemize}
    \item begins at 1,100
    \item increases by 120 bees per week in the first two weeks
    \item then increases by 180 bees per week until the goal is reached
\end{itemize}

According to this model, at the end of week \( w \) after the initial observation, where \( w > 2 \), which of the following functions gives the predicted number of bees still needed to reach the beekeeper’s goal?

\begin{enumerate}[A.]
    \item \( p(w) = 2,600 - 180w \)
    \item \( p(w) = 2,480 + 180w \)
    \item \( p(w) = 1,620 - 180w \)
    \item \( p(w) = -120 + 180w \)
\end{enumerate}
\end{frame}


\begin{frame}{\textbf{Answer Explanation: Bee Population Model }}
\small
\textbf{Step 1:} Initial population is 1,100.  
\textbf{Step 2:} In first 2 weeks, increase is \( 2 \times 120 = 240 \).  
So at the end of week 2, the population is:  
\[
1,100 + 240 = 1,340
\]

\pause

\textbf{Step 3:} For \( w > 2 \), increase is 180 bees per week.  
After \( w \) weeks, the population is:  
\[
\text{Population at week } w = 1,340 + 180(w - 2)
\]

\textbf{Step 4:} Number of bees still needed to reach 2,600 is:  
\[
p(w) = 2,600 - (1,340 + 180(w - 2))
\]

\pause

\textbf{Step 5:} Simplify the expression:
\[
p(w) = 2,600 - 1,340 - 180(w - 2)
= 1,260 - 180(w - 2)
\]
\[
= 1,260 - 180w + 360 = 1,620 - 180w
\]

\textbf{Correct answer: (C) \( \boxed{p(w) = 1,620 - 180w} \)}
\end{frame}


\begin{frame}{\textbf{Question: Discriminant and Inequality in Quadratics}}
\small
Given the quadratic equation:
\[
7rx^2 - 28sx + 7r
\]
In the equation:
\begin{itemize}
    \item \( r \) is a positive constant
    \item \( s \) is a negative constant
    \item The equation has two solutions
\end{itemize}

If \( \frac{r}{s} > k \), what is the minimum value of \( k \)?
\end{frame}


\begin{frame}
{\textbf{Answer Explanation: Discriminant and Inequality}}
\small
\textbf{Step 1:} A quadratic equation has two distinct real solutions if its discriminant is positive:
\[
\Delta = b^2 - 4ac > 0
\]

\textbf{Here:}
\[
a = 7r, \quad b = -28s, \quad c = 7r
\]

\pause

\[
(-28s)^2 - 4(7r)(7r) > 0
\]
\[
784s^2 - 196r^2 > 0
\]

Divide both sides by 196:
\[
4s^2 - r^2 > 0
\Rightarrow 4s^2 > r^2
\Rightarrow \frac{r^2}{s^2} < 4
\Rightarrow \left| \frac{r}{s} \right| < 2
\]

\[\frac{r}{s}>-2\]

\end{frame}


\begin{frame}{\textbf{Question: Map Scale Transformation}}
\small
A square map has a side length of 55 inches, and 1 inch on the map represents an actual distance of 13 miles.

A larger version of the same map is printed as a square with the side length 90\% longer than the side length of the previous map.

On the larger map, which of the following is closest to the actual distance, in miles, represented by 1 inch?

\begin{enumerate}[(A)]
    \item 1.30
    \item 6.84
    \item 24.70
    \item 49.50
\end{enumerate}
\end{frame}


\begin{frame}{\textbf{Answer Explanation: Map Scale Transformation}}
\textbf{Step 1:} Original scale is:
\[
\frac{13 \text{ miles}}{1 \text{ inch}} \text{ for a 55-inch map}
\]

\textbf{Step 2:} New map side is \( 55 \times 1.9 = 104.5 \) inches (90\% longer).

\pause

\textbf{Step 3:} Since the physical map is larger, each inch on the larger map covers fewer miles. We apply a scale factor:

\[
\text{Scale factor} = \frac{55}{104.5}
\Rightarrow \text{New miles per inch} = 13 \times \frac{55}{104.5} \approx 6.84
\]

\textbf{Answer:} \( \boxed{6.84} \)
\end{frame}


\begin{frame}{\textbf{Question: Isosceles Triangle and Congruent Angles}}
In isosceles triangle \( ABC \), sides \( AB \) and \( AC \) are congruent. Point \( D \) divides side \( BC \) such that the length of \( \overline{BD} \) is \( \frac{4}{7} \) of the length of \( \overline{BC} \).

Point \( E \) lies on side \( AB \) and point \( F \) lies on side \( AC \) such that when segments \( DE \) and \( DF \) are drawn, angle \( \angle BED \) is congruent to angle \( \angle CFD \).

If the length of \( \overline{BE} \) is 8, what is the length of \( \overline{CF} \)?

\begin{enumerate}[(A)]
    \item 6
    \item 14
    \item 32
    \item 56
\end{enumerate}
\end{frame}



\begin{frame}{\textbf{Answer Explanation: Isosceles Triangle and Congruent Angles}}
\textbf{Step 1:} Since \( \angle BED \cong \angle CFD \), and triangles \( \triangle BED \cong \triangle CFD \), then the triangles are similar by AA.

\pause

\textbf{Step 2:} We are told that \( BD : DC = 4 : 3 \), since \( BD = \frac{4}{7}BC \), meaning \( DC = \frac{3}{7}BC \). This is the same as the ratio \( \frac{BE}{CF} \).

\[
\frac{BE}{CF} = \frac{4}{3}
\Rightarrow \frac{8}{x} = \frac{4}{3} \Rightarrow x = \frac{8 \times 3}{4} = 6
\]

\textbf{Answer:} \( \boxed{6} \)
\end{frame}



\begin{frame}{\textbf{Question: Simplifying Radical Expressions}}

\[
\frac{\sqrt[3]{x^{17} y^5}}{5x^2 \left( \sqrt[8]{y^8} \right)}
\]

For all positive values of \( x \) and \( y \), the given expression is equivalent to which of the following?

\begin{enumerate}
    \item \( \frac{\left( x^{\frac{14}{3}} y^{\frac{5}{3}} \right)}{5xy} \)
    \item \( \frac{\sqrt[3]{x^{11} y^2}}{5} \)
\end{enumerate}

\begin{enumerate}[(A)]
    \item I only
    \item II only
    \item I and II
    \item Neither I nor II
\end{enumerate}

\end{frame}



\begin{frame}{Answer}


\textbf{Answer:} \( \boxed{C} \)
\end{frame}


\begin{frame}{Linear Equation from Table — Parameterized}
    The table shows three values of \(x\) and their corresponding values of \(y\), where \(s\) is a constant.

    \vspace{1em}
    \begin{center}
    \begin{tabular}{|c|c|}
        \hline
        \(x\) & \(y\) \\
        \hline
        \(-2s\) & 24 \\
        \(-s\) & 21 \\
        \(s\) & 15 \\
        \hline
    \end{tabular}
    \end{center}

    \vspace{1em}
    There is a linear relationship between \(x\) and \(y\).  
    Which of the following equations represents this relationship?

    \begin{itemize}
        \item[\textbf{A.}] \(sx + 3y = 18s\)
        \item[\textbf{B.}] \(3x + sy = 18s\)
        \item[\textbf{C.}] \(3x + sy = 18\)
        \item[\textbf{D.}] \(sx + 3y = 18\)
    \end{itemize}
\end{frame}


\begin{frame}{Answer: Linear Equation from Table}
    Use two known points:  
    \(x = -2s, y = 24\)  
    \(x = -s, y = 21\)

    \textbf{Step 1: Find the slope}

    \[
    m = \frac{21 - 24}{-s - (-2s)} = \frac{-3}{s}
    \]

    \textbf{Step 2: Use point-slope form with point \((-s, 21)\)}  
    \[
    y - 21 = \frac{-3}{s}(x + s)
    \]

    \textbf{Step 3: Rearrange to standard form}
    \[
    sy - 21s = -3x - 3s \quad \Rightarrow \quad 3x + sy = 18s
    \]

    \vspace{1em}
    \textbf{Correct Answer: \textcolor{novelPurple}{B. \(3x + sy = 18s\)}}
\end{frame}


\begin{frame}{Function Modeling: Carpet Cleaner Rental Cost}
    The cost of renting a carpet cleaner is \(\$52\) for the first day and \(\$26\) for each additional day.  
    Which of the following functions gives the cost \(C(d)\), in dollars, of renting the carpet cleaner for \(d\) days,  
    where \(d\) is a positive integer?

    \vspace{1em}
    \begin{itemize}
        \item[\textbf{A.}] \(C(d) = 26d + 26\)
        \item[\textbf{B.}] \(C(d) = 26d + 52\)
        \item[\textbf{C.}] \(C(d) = 52d - 26\)
        \item[\textbf{D.}] \(C(d) = 52d + 78\)
    \end{itemize}
\end{frame}


\begin{frame}{Answer: Carpet Cleaner Rental Function}
    \textbf{Step 1: Understand the scenario}

    \begin{itemize}
        \item The first day costs \(\$52\).
        \item Each additional day costs \(\$26\).
    \end{itemize}

    \textbf{Step 2: Break the cost down}

    \begin{itemize}
        \item If \(d = 1\), the cost is \(\$52\).
        \item If \(d = 2\), the cost is \(52 + 26 = 78\).
        \item For general \(d\), think of:
        \[
        C(d) = 52 + 26(d - 1)
        \]
        \item Simplify:
        \[
        C(d) = 52 + 26d - 26 = 26d + 26
        \]
    \end{itemize}

    \textbf{Correct Answer: \textcolor{novelPurple}{A. \(C(d) = 26d + 26\)}}
\end{frame}

\begin{frame}{Unit Conversion: Acceleration}
    The speed of a vehicle is increasing at a rate of \(7.3\) meters per second squared.  
    What is this rate, in \textbf{miles per minute squared}, rounded to the nearest tenth?

    \vspace{0.5em}
    (Use \(1~\text{mile} = 1{,}609~\text{meters}\))

    \vspace{1em}
    \begin{itemize}
        \item[\textbf{A.}] \(0.3\)
        \item[\textbf{B.}] \(16.3\)
        \item[\textbf{C.}] \(195.8\)
        \item[\textbf{D.}] \(220.4\)
    \end{itemize}
\end{frame}


\begin{frame}{Answer: Unit Conversion}
    \textbf{Given:}  
    \[
    7.3~\text{m/s}^2
    \]
    Convert to miles per minute squared:

    \vspace{0.5em}
    \[
    \frac{7.3~\text{m}}{1~\text{s}^2} \times \frac{1~\text{mile}}{1{,}609~\text{m}} \times \left(60~\text{s}\right)^2
    \]

    \vspace{0.5em}
    \[
    = \frac{7.3}{1{,}609} \times 3{,}600 = 16.328 \approx \boxed{16.3}
    \]

    \textbf{Correct Answer: \textcolor{novelPurple}{B. \(16.3\)}}
\end{frame}







\end{document}